\section{Fase 4. Desarrollo del Director del juego}
\label{sec:Fase4}

Hasta ahora disponemos de un director del juego básico, que lanza los dos motores necesarios y poco más. Ahora lo ajustaremos para incluir todas sus tareas. Esta fase saca del código todos los parámetros que gobiernan el ritmo de la partida y los convierte en propiedades del cartucho, editables desde el diseñador gráfico. La motivación es doble: por un lado, el simulador de PC y el futuro sistema sobre ESP32 comparten toda la lógica del motor de estados, que es determinista en \textit{ticks}, pero sus procesadores trabajan a velocidades muy distintas y pueden exigir pausas de reloj diferentes para ofrecer el mismo \textit{feel} de juego; por otro, la fase siguiente (el generador del cartucho compacto) necesita que \textbf{todo} lo configurable viaje en el GWI, incluido lo que hasta ahora solo vivía en el setup o en constantes del código.

Cerraremos en esta fase dos temas: 
\begin{itemize}
	\item Esta fase ajusta el Director del juego: la velocidad por plataforma, la aceleración, la política de rutas y la combinación de pulsaciones dejan de ser constantes del código o del setup y pasan a ser propiedades del cartucho, editables en el diseñador y visibles para el motor al instante, sin guardar ni recargar. 
	
	\item Como objetivo secundario, cerraremos todos los temas pendientes para la generación del cartucho intermedio, que debe llevar toda la información necesaria para, en la siguiente fase, generar el cartucho compacto.
\end{itemize}

Con esta fase los parámetros salen definitivamente del setup y del código: nacen en
el cartucho (GWI) y viajan al Bloque 1 del GWC (F5-I02); lo que el GWI no declara
viaja como nulo (F5-I13) y el valor por defecto (los 500 ms históricos, una sola ruta...) lo aplica quien lee, igual en el simulador y en el ESP32. La fuente de verdad
durante la autoría es siempre el diseñador gráfico: cualquier edición en la ficha del
cartucho se sincroniza con el motor al instante, sin guardar ni recargar.


\clearpage
